	\section*{Введение}
\addcontentsline{toc}{section}{Введение}

\subsection*{Для кого это пособие}

Это пособие написано для студентов инженерных специальностей и исследователей,
которые хотят научиться использовать Python для обработки экспериментальных
данных. Мы не будем изучать Python «в вакууме» — мы будем решать конкретную,
осязаемую физическую задачу, и по ходу дела освоим всё необходимое.

Предполагается, что ты:
\begin{itemize}
	\item имеешь базовое представление о программировании (знаешь, что такое
	переменная, цикл, функция);
	\item умеешь работать с файлами в своей операционной системе;
	\item не боишься открыть командную строку.
\end{itemize}

Если ты никогда не программировал на Python — не страшно. Мы будем двигаться плавно, 
шаг за шагом, объясняя каждую строчку кода и её физический смысл.

\subsection*{Что мы будем делать}

Мы создадим программу для анализа данных с двумя проволочными поляризаторами —
устройствами, которые используются в терагерцовой (ТГц) спектроскопии для управления
поляризацией излучения (принципиальная схема эксперимента представлена на \figref{fig:0_1}).

\smartfigure{fig0.1.png}{Принципиальная схема эксперимента с двумя проволочными поляризаторами на пути ТГц излучения.}{fig:0_1}

Поляризаторы вращаются друг относительно друга, пропуская разную долю излучения.
Наша конечная цель — проверить, насколько реальное пропускание системы из двух поляризаторов
соответствует теоретической модели Бланко (Blanco model).

\textbf{Исходные данные:} текстовые файлы с временны́ми сигналами.
В каждом файле — два столбца: время (в пикосекундах) и амплитуда сигнала.
Имена файлов содержат информацию об углах поворота поляризаторов,
номере измерения и номере фонового сигнала.

\textbf{Что должна делать программа:}
\begin{enumerate}
	\item Загружать все файлы из папки эксперимента.
	\item Извлекать из имён файлов информацию об углах поворота.
	\item Усреднять повторные измерения для уменьшения шума.
	\item Удалять постоянную составляющую сигнала (DC bias).
	\item Применять оконную фильтрацию для отсечения ложных отражений.
	\item Выполнять быстрое преобразование Фурье (БПФ) для перехода в частотную область.
	\item Вычислять спектры пропускания.
	\item Строить графики спектров и угловой зависимости.
	\item Сравнивать эксперимент с теорией Бланко и подбирать параметры модели.
\end{enumerate}

\subsection*{Что ты изучишь по ходу дела}

В процессе работы ты освоишь:
\begin{itemize}
	\item работу с файловой системой в Python (модуль \texttt{pathlib});
	\item парсинг (разбор) текстовых данных и регулярные выражения;
	\item векторизованные вычисления с массивами \textbf{NumPy};
	\item основы цифровой обработки сигналов: БПФ, оконные функции;
	\item построение профессиональных научных графиков с \textbf{Matplotlib};
	\item численную оптимизацию и аппроксимацию с помощью \textbf{SciPy};
	\item организацию чистого и расширяемого кода.
\end{itemize}

\subsection*{Структура пособия}

Пособие построено по принципу «от простого к сложному». Каждый раздел —
это законченный этап работы, после которого у тебя будет работающий скрипт
с новыми возможностями. Сначала мы разбираем физический или математический смысл шага, затем пишем код, выполняем его и анализируем результат.

\begin{enumerate}
	\item \textbf{Подготовка инструментов} — устанавливаем Python и библиотеки.
	\item \textbf{Структура проекта} — создаём файлы и настраиваем конфигурацию.
	\item \textbf{Загрузка данных} — учимся читать файлы и парсить имена.
	\item \textbf{Первичная визуализация} — смотрим, как выглядят сигналы во времени.
	\item \textbf{Обработка сигналов} — центрирование, окна, БПФ и расчет спектров.
	\item \textbf{Теоретическая модель} — реализуем формулы модели Бланко.
	\item \textbf{Угловая зависимость} — строим зависимости и сравниваем с теорией.
	\item \textbf{Оптимизация} — подбираем физические параметры автоматически.
	\item \textbf{Финальная сборка} — объединяем весь код в удобное приложение.
\end{enumerate}

\subsection*{Как работать с пособием}

\warning{\textbf{Подготовка материалов:} К данной методичке прилагается архив с исходными кодами и тестовыми данными (папки \texttt{src} и \texttt{data}). Перед началом работы \textbf{обязательно скачай этот архив и распакуй его} в удобное место (например, на Рабочий стол). В тексте мы часто будем ссылаться на исходные файлы из этого архива.}

\begin{enumerate}
	\item Читайте текст последовательно, не перескакивая через разделы.
	\item Пиши код самостоятельно или копируй из листингов в свой проект.
	\item Запускай скрипты и проверяй, совпадает ли твой результат с картинками в пособии.
	\item Экспериментируй — меняй параметры фильтров или графиков, смотри, как реагирует программа.
	\item Если что-то сломалось — не паникуй, читай текст ошибок в консоли, ищи блоки \textbf{Важно} в тексте.
\end{enumerate}

\textbf{Важное правило:} код в пособии — это не «священный текст», который
нужно бездумно копировать. Это инструмент, который ты должен понять и
научиться использовать. Поэтому каждая важная строчка будет подробно разобрана.

\subsection*{Обозначения}

В тексте используются следующие визуальные маркеры:
\begin{itemize}
	\item \texttt{моноширинный шрифт} — имена файлов, переменные, функции;
	\item \textbf{жирный шрифт} — важные термины и определения;
	\item \note{Примечание} — интересная физическая или программистская справка;
	\item \warning{Важно} — подводные камни, частые ошибки или критические требования;
	\item \tip{Совет} — лайфхак или более элегантный способ сделать то же самое.
\end{itemize}

